% tex/sections/outlook.tex
% ──────────────────────────────────────────────────────────────────────────────
\chapter{Outlook}
\label{chap:outlook}
% ──────────────────────────────────────────────────────────────────────────────
%
% TODO: Write outlook.
%
% Suggested structure:
%   - Future work directly motivated by this thesis
%   - Broader research directions opened up
%   - Proposed experiments or improvements to methodology

Based on the results and methods presented in this thesis, there are many open lines of inquiry. Because the model described is a basic two layer model, a next step could be to add multiple layers, especially because cracks in snow usually propagate in a weak layer \parencite{schweizerSnowAvalancheFormation2003}. \par

There is also the possibility of adding topography or a slope angle to the simulation, because as indicated by \cite{anceySnowAvalanches2001,trottetTransitionSubRayleighAnticrack2022}, avalanches are gravity driven and the crack propagation regimes can change greatly with the presence of topography. \par 

There is the possibility for transition between the different propagation regimes, meaning that the propagation speed changes, as described by \cite{trottetTransitionSubRayleighAnticrack2022}, which is not only driven by topography and slope angle. Another next step would be to have source propagation accelerate with increasing propagation distance, instead of having a constant propagation speed, which is not very plausible on a mountain slope. \par 
